\chapter{Boyer--Moore}

\scan{9}
L'algoritmo di Boyer--Moore (1977) risolve lo stesso problema di Knuth--Morris--Pratt
--- trovare tutte le occorrenze di un pattern $P$ dentro un testo $T$ --- ma cambia la
direzione dei confronti all'interno di un singolo allineamento. Da questa sola idea
discendono due euristiche di riposizionamento che, nei casi favorevoli, permettono di
non leggere affatto alcuni caratteri del testo. Per tutto il capitolo poniamo
\[
  m=\abs{T},\qquad n=\abs{P},
\]
con indicizzazione a partire da $1$; $\alfa$ è l'alfabeto e $P^r$ denota il pattern
\emph{rovesciato}.

\section{L'idea: confronti da destra a sinistra}

Il testo $T$ viene percorso da sinistra a destra, cioè gli allineamenti di $P$ su $T$
avanzano come nell'algoritmo naive e in KMP; \emph{dentro} un allineamento, però, il
pattern $P$ viene confrontato con il testo \emph{da destra a sinistra}. È questa
l'idea che si aggiunge rispetto agli algoritmi visti finora.

\begin{figure}[H]
\centering
% Un passo di Boyer--Moore: confronti da destra a sinistra, mismatch x/y, shift.
\begin{tikzpicture}[x=0.38cm,y=0.38cm,
    testo/.style   = {line width=0.9pt},
    patt/.style    = {line width=0.9pt,deepsea},
    taccaT/.style  = {line width=0.8pt},
    taccaP/.style  = {line width=0.8pt,deepsea},
    forte/.style   = {line width=1.8pt,deepsea},
    et/.style      = {font=\footnotesize},
    ets/.style     = {font=\scriptsize,graydark}]

  %% ------------------------------------------------- riga 1: il testo T
  \draw[testo] (0,0) -- (25,0);
  \node[et,anchor=west] at (25.5,0) {$T$};
  \draw[taccaT] (11,-0.7) -- (11,0.7);
  \draw[taccaT] (12,-0.7) -- (12,0.7);
  \draw[taccaT] (15,-0.7) -- (15,0.7);
  \node[et,anchor=south] at (11.5,0.75) {$x$};
  \node[et,anchor=south] at (13.5,0.75) {$\alpha$};

  %% ------------------------- direzione dei confronti (da destra a sinistra)
  \draw[graydark,line width=0.6pt,-{Latex[length=4pt]}] (15,-2.2) -- (11,-2.2);
  \node[ets,anchor=south] at (13,-2.1) {confronti};

  %% ---------------------------- riga 2: il pattern P, allineamento corrente
  \draw[patt] (6,-4.6) -- (15,-4.6);
  \draw[taccaP] (6,-5.25) -- (6,-3.95);
  \draw[taccaP] (11,-5.25) -- (11,-3.95);
  \draw[taccaP] (12,-5.25) -- (12,-3.95);
  \draw[forte]  (15,-5.5) -- (15,-3.7);
  \node[et,anchor=south] at (11.5,-3.9) {$y$};
  \node[et,anchor=south] at (13.5,-3.9) {$\alpha$};
  \node[et,anchor=north] at (15,-5.75) {$P$};

  %% --------------------------------- porzione di P non ancora confrontata
  \draw[graydark,line width=0.6pt]
        (6.8,-5.9) .. controls (8.2,-6.8) and (9.6,-5.2) .. (11,-5.9);
  \node[ets,anchor=north] at (8.9,-6.9) {non ancora confrontata};

  %% ------------------------------- riga 2: il pattern P dopo lo shift
  \draw[patt] (15.8,-4.6) -- (24.8,-4.6);
  \draw[taccaP] (15.8,-5.25) -- (15.8,-3.95);
  \draw[forte]  (24.8,-5.5) -- (24.8,-3.7);
  \node[et,anchor=north] at (24.8,-5.75) {$P$};

  %% ------------------------------------------------- annotazione a destra
  \node[et,anchor=west,align=left] at (26.1,-2.4)
       {ci sono casi\\[2pt] $\Oh(m/n)$};
\end{tikzpicture}

\caption{Un passo di Boyer--Moore: i confronti procedono da destra a sinistra,
$\alpha$ è il suffisso di $P$ già riconosciuto e il mismatch è fra il carattere
$y$ del pattern e il corrispondente carattere $x$ del testo; il pattern viene poi
fatto scorrere verso destra, nel disegno addirittura di tutta la propria lunghezza.}
\end{figure}

Eseguo dunque i confronti da destra e, quando trovo un mismatch, applico la
\emph{più conveniente} fra le due regole della sezione seguente, cioè quella che
produce il \emph{massimo shift}.

\begin{osservazione}
Poiché dopo un mismatch il pattern può scorrere anche di tutta la propria lunghezza,
ci sono casi in cui i confronti eseguiti --- e quindi i caratteri di $T$
effettivamente letti --- sono soltanto $\Th(m/n)$: l'algoritmo può essere
\emph{sublineare} nella lunghezza del testo, cioè terminare senza aver mai letto
alcuni caratteri di $T$. Il preprocessing del pattern, $\Th(n)$, resta ovviamente
da pagare.
\end{osservazione}

\section{Le due regole}\label{sec:bm-regole}

\scan{10}
Fissiamo la notazione di un allineamento. Supponiamo che $P$ sia allineato con la
finestra di testo che termina in $T[k]$ e che i confronti da destra abbiano già
riconosciuto il suffisso
\[
  \alpha=\sub{P}{i}{n}
\]
(con $\alpha=\vuota$ quando $i=n+1$, cioè quando il primissimo confronto fallisce).
Il primo confronto che fallisce è allora quello fra il carattere $y=P[i-1]$ del
pattern e il carattere $x$ del testo che gli sta di fronte, con $x\neq y$: in tutto
il capitolo $x$ è dunque il carattere «cattivo» del \emph{testo} --- i due ruoli
sono scambiati rispetto al capitolo precedente, dove $x$ era il carattere del
pattern. Al mismatch il pattern viene spostato a destra applicando la più
conveniente fra le due regole seguenti.

\begin{definizione}[Bad character rule (BCR)]\label{def:bcr}
Dopo un mismatch con il carattere $x$ di $T$, si riposiziona il pattern cercando in
$P$ l'occorrenza di $x$ \emph{più a destra fra quelle a sinistra} del punto di
mismatch.
\end{definizione}

\begin{definizione}[Good suffix rule (GSR)]\label{def:gsr}
Dopo un mismatch che segue il suffisso $\alpha$ già riconosciuto in $P$, si
riposiziona $P$ cercando l'occorrenza di $\alpha$ più a destra che non sia un
suffisso di $P$.
\end{definizione}

Le due regole hanno costi di preprocessing molto diversi: la BCR si implementa
direttamente, mentre il preprocessing della GSR si riduce all'algoritmo $Z$
(\autoref{sec:bm-gsr}).

\begin{osservazione}[Di quanto si sposta il pattern]
Per ogni posizione $j$ e ogni $x\in\alfa$ poniamo
\[
  R_j(x) \defeq \max\set{\, p \;:\; 1\le p<j \ \text{e}\ P[p]=x \,},
  \qquad R_j(x)=0 \ \text{se un tale } p \text{ non esiste}.
\]
Con la notazione fissata sopra, la BCR sposta $P$ di $(i-1)-R_{i-1}(x)$ posizioni,
mentre la GSR lo sposta di $n-j_\alpha$, dove $j_\alpha$ è la posizione finale
dell'occorrenza di $\alpha$ individuata dalla regola. Nella \autoref{sec:bm-gsr}
rafforzeremo la regola, chiedendo in più che il carattere che precede
l'occorrenza interna sia diverso da quello che ha provocato il mismatch, e vedremo
che allora $j_\alpha$ è esattamente il valore $L'(i)$.
Quando una tale occorrenza \emph{non} esiste --- cioè quando $L'(i)=0$ --- lo
spostamento sarebbe di $n$ posizioni intere, e così com'è non è \emph{safe}: su
$P=\str{aca}$ e $T=\str{bcaca}$ farebbe saltare l'occorrenza in posizione~3. Il
caso va trattato a parte, come si vedrà nella \autoref{sec:bm-ricerca}.
Con questa avvertenza entrambi gli spostamenti sono $\ge 1$, e si applica sempre
il massimo dei due,
\[
  \mathit{shift} \;=\; \max\set{\,\mathit{shift}_{\mathrm{BCR}},\ \mathit{shift}_{\mathrm{GSR}}\,},
\]
il che è lecito perché entrambe le regole sono \emph{safe} --- nessuna occorrenza
viene persa --- e quindi lo è anche il loro massimo.
\end{osservazione}

\begin{esempio*}[Comportamento sublineare]
Siano $T=a^{m}$ e $P=b^{n}$. Ogni allineamento fallisce al primo confronto e il
carattere cattivo $a$ non compare affatto in $P$: la BCR dà $R_n(a)=0$ e fa scorrere
il pattern di tutta la sua lunghezza. Si esaminano così $m/n$ allineamenti con un
solo confronto ciascuno, cioè $\Th(m/n)$ caratteri di $T$ in tutto.
\end{esempio*}

\begin{osservazione}
All'estremo opposto, se $T=a^{m}$ e $P=a^{n}$ ogni allineamento produce
un'occorrenza completa e nessuna delle due regole permette di avanzare di più di una
posizione: l'algoritmo degenera nell'algoritmo naive, con costo quadratico
$\Th(mn)$.
\end{osservazione}

\textbf{N.B.} È possibile, aggiungendo una struttura dati, rendere l'algoritmo
lineare nel caso peggiore (Apostolico--Giancarlo) --- \emph{da approfondire}.

\section{Implementazione della BCR}

Il preprocessing riguarda il solo pattern $P$.

\begin{figure}[H]
\centering
% Preprocessing della BCR: da j si guarda all'indietro dentro P cercando x.
\begin{tikzpicture}[font=\footnotesize,
    tacca/.style = {line width=1pt},
    ets/.style   = {font=\scriptsize,graydark}]

  %% il pattern P
  \draw[deepsea,line width=2.2pt,-{Latex[length=5pt,width=4.5pt]}] (0,0) -- (10,0);
  \node[anchor=west] at (10.3,0) {$P$};
  \node[ets,anchor=north] at (0,-0.45) {$1$};
  \node[ets,anchor=north] at (9.8,-0.45) {$n$};

  %% le due posizioni: p (occorrenza trovata) e j (mismatch corrente)
  \draw[tacca] (6.4,-0.32) -- (6.4,0.32);
  \draw[tacca] (7.6,-0.32) -- (7.6,0.32);
  \node[anchor=north] at (6.4,-0.36) {$p$};
  \node[anchor=south] at (7.6,0.36) {$j$};
  \node[anchor=south] at (6.8,0.3) {$x$};

  %% si guarda all'indietro, da j verso p
  \draw[gray,dashed,line width=0.7pt,-{Latex[length=4.5pt]}]
        (7.6,-0.9) .. controls (7.6,-1.5) and (6.4,-1.5) .. (6.4,-0.9);
  \node[ets,anchor=north] at (7.0,-1.45) {$j-R_j(x)$};
\end{tikzpicture}

\caption{Preprocessing della BCR: dalla posizione $j$ si guarda all'indietro dentro
$P$ per individuare la posizione $p<j$ più a destra con $P[p]=x$.}
\end{figure}

Per ogni $j\in\set{1,\dots,n}$ e per ogni $x\in\alfa$ serve il massimo indice $p$
tale che $p<j$ e $P[p]=x$, cioè esattamente il valore $R_j(x)$ definito sopra. Ci
sono due modi di memorizzarlo.

\begin{itemize}
  \item \textbf{Matrice.} Si tabula direttamente $R_j(x)$ in una matrice
  $n\times\abs{\alfa}$: costa $\Th(n\,\abs{\alfa})$ in spazio e in tempo di
  preprocessing, ma risponde a ogni interrogazione in tempo costante.

  \item \textbf{Liste.} Per ogni carattere $x$ che compare in $P$ si mantiene la
  lista delle posizioni di $P$ in cui $x$ occorre, in ordine decrescente. Lo spazio
  scende a $\Th(n+\abs{\alfa})$, ma per calcolare $R_j(x)$ bisogna scorrere la
  lista.
\end{itemize}

La seconda soluzione sembra quadratica; facendo bene i conti, però, con un'analisi a
complessità ammortizzata si vede che in realtà è lineare.

\begin{osservazione}[Perché il costo ammortizzato è lineare]
Poiché la lista di $x$ è ordinata per posizione decrescente, per calcolare $R_j(x)$
basta scorrerla dalla testa scartando le posizioni $\ge j$. Se il mismatch cade in
$P[j]$, in quel momento sono già stati eseguiti $n-j+1$ confronti (le posizioni
$n,\dots,j$ di $P$), e le posizioni scartate sono tutte comprese fra $j$ e $n$:
sono quindi al più $n-j+1$. Il costo della scansione della lista è perciò
proporzionale al numero di confronti appena eseguiti, e sull'intera esecuzione il
sovraccosto complessivo resta lineare nel numero totale di confronti:
asintoticamente le due implementazioni si equivalgono.
\end{osservazione}

\section{Implementazione della GSR}\label{sec:bm-gsr}

\scan{11}
Resta da capire quale sia il preprocessing sufficiente a implementare la
\emph{good suffix rule}. La risposta è: l'algoritmo $Z$ applicato a $P^r$.

\begin{definizione}[$L'(i)$]\label{def:bm-lprimo}
Siano $P\in\alfa^*$ e $i\in\set{2,\dots,n}$. Allora $L'(i)$ è il \emph{massimo}
indice $j$ tale che
\begin{enumerate}
  \item $\sub{P}{i}{n}$ è un suffisso di $\sub{P}{1}{j}$;
  \item $P[i-1]\neq P[\,j-(n-i)-1\,]$;
\end{enumerate}
e $L'(i)=0$ se nessun indice soddisfa 1) e 2).
\end{definizione}

\begin{osservazione}
La sola condizione~1), con il vincolo $j\le n-1$, definirebbe $L(i)$, cioè la
versione «debole» della regola: l'occorrenza interna più a destra del suffisso già
riconosciuto (il vincolo è indispensabile: $\sub{P}{i}{n}$ è sempre suffisso di
$P$, quindi senza di esso il massimo sarebbe banalmente $j=n$). È la condizione~2)
a giustificare l'apice in $L'(i)$: essa impone che il carattere che precede
l'occorrenza interna sia \emph{diverso} da $P[i-1]$, cioè dal carattere che
nell'allineamento corrente ha provocato il mismatch; senza questo vincolo lo
spostamento riproporrebbe lo stesso mismatch e sarebbe quindi sprecato.

Gli indici della condizione~2) si leggono così: se $\sub{P}{i}{n}$, che ha lunghezza
$n-i+1$, è un suffisso di $\sub{P}{1}{j}$, allora occupa le posizioni
$j-(n-i),\dots,j$, e il carattere che la precede sta in posizione $j-(n-i)-1$.

Restano sottintese due convenzioni: se $j-(n-i)-1=0$ la posizione non esiste e la
condizione~2) si intende soddisfatta; $L'(i)$ non è definito per $i=1$, perché la
condizione~2) chiamerebbe in causa $P[0]$. Si noti infine che la~2) esclude
automaticamente il valore $j=n$: in tal caso diventerebbe $P[i-1]\neq P[i-1]$, che è
falsa. Dunque $L'(i)\le n-1$ senza bisogno di imporlo, ed è per questo che lo
spostamento $n-L'(i)$ della GSR è sempre positivo.
\end{osservazione}

\begin{figure}[H]
\centering
% Lettura grafica di L'(i): la stessa stringa P[i..n] occorre due volte in P
% (come suffisso e come occorrenza interna che termina in L'(i)) e i due
% caratteri che la precedono sono diversi.
\begin{tikzpicture}[x=1.3cm,font=\footnotesize,
    tacca/.style = {line width=1pt},
    barra/.style = {line width=2.4pt,graydark},
    ets/.style   = {font=\scriptsize,graydark}]

  %% il pattern P: posizione 1 in x=0, posizione n in x=10.4
  \draw[deepsea,line width=2.2pt,-{Latex[length=5pt,width=4.5pt]}] (0,0) -- (10.6,0);
  \node[anchor=west] at (10.9,0) {$P$};

  %% tacche
  \foreach \t in {2.4,3.6,6.0,6.8,8.0} \draw[tacca] (\t,-0.3) -- (\t,0.3);

  %% i due caratteri che precedono le due occorrenze
  \node[anchor=south] at (3.0,0.16) {$z$};
  \node[anchor=south] at (7.4,0.16) {$y$};

  %% le due occorrenze della stessa stringa P[i..n], di uguale lunghezza
  \draw[barra] (3.6,0.55) -- (6.0,0.55);
  \draw[tacca,graydark] (3.6,0.38) -- (3.6,0.72);
  \draw[tacca,graydark] (6.0,0.38) -- (6.0,0.72);
  \draw[barra] (8.0,0.55) -- (10.4,0.55);
  \draw[tacca,graydark] (8.0,0.38) -- (8.0,0.72);
  \draw[tacca,graydark] (10.4,0.38) -- (10.4,0.72);
  \node[anchor=south] at (4.8,0.68) {$\sub{P}{i}{n}$};
  \node[anchor=south] at (9.2,0.68) {$\sub{P}{i}{n}$};

  %% condizione 2): i caratteri che precedono le due occorrenze sono diversi
  \draw[gray,dashed,line width=0.7pt]
        (3.0,0.6) .. controls (3.9,2.6) and (6.5,2.6) .. (7.4,0.6);
  \node[fill=white,inner sep=3pt] at (5.2,2.1) {$\neq$};

  %% etichette delle posizioni
  \node[ets,anchor=north] at (3.0,-0.34) {$L'(i)-(n-i)-1$};
  \node[anchor=north] at (6.0,-0.34) {$L'(i)$};
  \node[anchor=north] at (7.4,-0.34) {$i-1$};
  \node[anchor=north west] at (8.0,-0.34) {$i$};

  %% condizione 1): P[i..n] e' un suffisso di P[1..L'(i)]
  \draw[decorate,decoration={brace,mirror,amplitude=5pt},gray,line width=0.7pt]
        (0,-1.05) -- (6.0,-1.05);
  \node[ets,anchor=north] at (3.0,-1.32) {$\sub{P}{1}{L'(i)}$};
\end{tikzpicture}

\caption{Lettura grafica della definizione di $L'(i)$: la stessa stringa
$\sub{P}{i}{n}$ compare due volte in $P$ --- come suffisso e come occorrenza interna
che termina in $L'(i)$ --- e i due caratteri che la precedono sono diversi.}
\end{figure}

\begin{definizione}[{{$N_j[P]$}}]\label{def:bm-nj}
$N_j[P]$ è la lunghezza del massimo suffisso di $\sub{P}{1}{j}$ che è anche un
suffisso di $P$, ossia
\[
  N_j[P]=\max\set{\ell\ge 0 \;:\; \sub{P}{j-\ell+1}{j}=\sub{P}{n-\ell+1}{n}},
  \qquad j=1,\dots,n .
\]
\end{definizione}

I valori $N_j$ sono l'analogo «riflesso» dei valori $Z$: $\Zv{k}$ guarda i prefissi
da sinistra, $N_j$ guarda i suffissi da destra. In particolare $N_n[P]=n$.

\begin{figure}[H]
\centering
% N_j[P] = |alpha|, con alpha suffisso di P[1..j] e insieme suffisso di P.
\begin{tikzpicture}[font=\footnotesize,
    tacca/.style = {line width=1pt},
    ets/.style   = {font=\scriptsize,graydark}]

  %% il pattern P: posizione 1 in x=0, posizione n in x=10.4
  \draw[deepsea,line width=2.2pt,-{Latex[length=5pt,width=4.5pt]}] (0,0) -- (10.8,0);
  \node[anchor=west] at (11.1,0) {$P$};

  %% tacche agli estremi delle due occorrenze di alpha
  \foreach \t in {2.3,4.3,8.4,10.4} \draw[tacca] (\t,-0.3) -- (\t,0.3);

  %% occorrenza che termina in j
  \draw[decorate,decoration={brace,amplitude=5pt,raise=3pt},graydark,line width=0.7pt]
        (2.3,0) -- (4.3,0);
  \node[anchor=south] at (3.3,0.36) {$\alpha$};

  %% occorrenza che termina in n: alpha e' anche un suffisso di P
  \draw[decorate,decoration={brace,amplitude=5pt,raise=3pt},graydark,line width=0.7pt]
        (8.4,0) -- (10.4,0);
  \node[anchor=south] at (9.4,0.36) {$\alpha$};

  %% le due occorrenze sono la stessa stringa
  \draw[gray,dashed,line width=0.7pt]
        (3.3,0.78) .. controls (4.4,1.9) and (8.3,1.9) .. (9.4,0.78);
  \node[fill=white,inner sep=3pt] at (6.35,1.62) {$=$};

  %% etichette delle posizioni
  \node[ets,anchor=north] at (0,-0.3) {$1$};
  \node[anchor=north] at (4.3,-0.3) {$j$};
  \node[ets,anchor=north] at (10.4,-0.3) {$n$};
\end{tikzpicture}

\caption{$N_j[P]=\abs{\alpha}$, dove $\alpha$ è la più lunga stringa che è
simultaneamente suffisso di $\sub{P}{1}{j}$ e suffisso di $P$.}
\end{figure}

\begin{teorema}\label{thm:bm-lprimo-nj}
Per ogni $i\in\set{2,\dots,n}$, $L'(i)$ è il massimo $j$ tale che
\[
  N_j[P]=\abs{\sub{P}{i}{n}}=n-i+1 .
\]
\end{teorema}

\begin{dimostrazione}
Basta mostrare che, per ogni indice $j$, le condizioni 1) e 2) della
Definizione~\ref{def:bm-lprimo} valgono per quel $j$ se e solo se $N_j[P]=n-i+1$.

Se valgono 1) e 2), la 1) dice che $\sub{P}{j-(n-i)}{j}=\sub{P}{i}{n}$, dunque
$\sub{P}{1}{j}$ ha un suffisso di lunghezza $n-i+1$ che è anche suffisso di $P$ e
quindi $N_j[P]\ge n-i+1$; se fosse $N_j[P]\ge n-i+2$ il match si estenderebbe di un
carattere verso sinistra, cioè $P[j-(n-i)-1]=P[i-1]$, contro la 2).

Viceversa, se $N_j[P]=n-i+1$ allora il suffisso comune di quella lunghezza è
esattamente $\sub{P}{j-(n-i)}{j}=\sub{P}{i}{n}$, cioè vale la 1); e la
\emph{massimalità} di $N_j$ impedisce l'estensione a sinistra, cioè
$P[j-(n-i)-1]\neq P[i-1]$, che è la 2) (se $j-(n-i)-1=0$ la condizione è soddisfatta
per convenzione, perché il match non può estendersi oltre l'inizio di $P$).

Le due condizioni sono equivalenti indice per indice, quindi hanno lo stesso massimo.
Si noti che per $i\ge 2$ il valore $j=n$ è automaticamente escluso, poiché
$N_n[P]=n>n-i+1$.
\end{dimostrazione}

Da qui, e dalla definizione stessa di $N_j$, segue la riduzione all'algoritmo $Z$:
per ogni $1\le j\le n-1$,
\[
  \boxed{\;N_j[P]=\Zv{n-j+1}[P^r]\;}
\]
Il solo valore escluso è $j=n$: chiamerebbe in causa $\Zv{1}[P^r]$, che la
definizione della funzione $Z$ non contempla. Non è una perdita, perché
$N_n[P]=n$ è noto a priori e nel seguito $N_j$ serve soltanto per $j\le n-1$.

\begin{osservazione}
Rovesciando $P$, la posizione $j$ di $P$ diventa la posizione $n-j+1$ di $P^r$; un
suffisso di $\sub{P}{1}{j}$ diventa un prefisso di $\sub{P^r}{n-j+1}{n}$ e un
suffisso di $P$ diventa un prefisso di $P^r$. Quindi il più lungo suffisso di
$\sub{P}{1}{j}$ che è anche suffisso di $P$ ha la stessa lunghezza del più lungo
prefisso di $\sub{P^r}{n-j+1}{n}$ che è anche prefisso di $P^r$, cioè
$\Zv{n-j+1}[P^r]$.

Poiché l'algoritmo $Z$ su una stringa di lunghezza $n$ costa $\Th(n)$, tutto il
preprocessing della GSR costa $\Th(n)$ in tempo e in spazio.
\end{osservazione}

\section{Lo spostamento della GSR e il ciclo di ricerca}\label{sec:bm-ricerca}

Sapere calcolare $L'(i)$ non basta ancora a far girare l'algoritmo: manca il caso in
cui $L'(i)=0$, cioè quando il suffisso $\alpha$ già riconosciuto non ha alcuna
ulteriore occorrenza utile dentro $P$. In quel caso non si può buttare via tutto:
bisogna allineare il più lungo prefisso di $P$ che sia anche suffisso di $\alpha$.

\begin{definizione}[$\ell'(i)$]\label{def:bm-lminuscolo}
Per $i\in\set{2,\dots,n+1}$, $\ell'(i)$ è la lunghezza del più lungo suffisso di
$\sub{P}{i}{n}$ che è anche un prefisso di $P$, e $\ell'(i)=0$ se un tale suffisso
non esiste. In particolare $\ell'(i)\le n-i+1$.
\end{definizione}

Anche $L'$ e $\ell'$ si leggono direttamente dal vettore $N$:
\begin{align*}
  L'(i) &= \max\set{\,j \;:\; 1\le j\le n-1,\ N_j[P]=n-i+1\,},\\
  \ell'(i) &= \max\set{\,j \;:\; 1\le j\le n-i+1,\ N_j[P]=j\,},
\end{align*}
dove il massimo su un insieme vuoto vale $0$; entrambi i vettori si riempiono con
una singola scansione di $N$, quindi in tempo $\Th(n)$. La seconda formula si
giustifica così: $N_j[P]=j$ significa che $\pre{P}{j}$ è anche un suffisso di $P$, e
la limitazione $j\le n-i+1$ impone che quel suffisso stia dentro $\sub{P}{i}{n}$.

Lo spostamento prodotto dalla GSR è allora
\[
  \sigma(i)\defeq
  \begin{cases}
    n-\ell'(2) & \text{se } i=1 \quad (\text{occorrenza completa}),\\[3pt]
    n-L'(i)  & \text{se } 2\le i\le n \ \text{e}\ L'(i)>0,\\[3pt]
    n-\ell'(i) & \text{se } 2\le i\le n \ \text{e}\ L'(i)=0,\\[3pt]
    1        & \text{se } i=n+1 \quad (\alpha=\vuota).
  \end{cases}
\]
L'ultimo caso è obbligato: se il primo confronto fallisce, nessun suffisso di $P$ è
stato riconosciuto e la GSR non porta alcuna informazione. Negli altri casi
$\sigma(i)\ge 1$, perché $L'(i)\le n-1$ e $\ell'(i)\le n-i+1\le n-1$.

\begin{algorithm}[H]
\caption{Boyer--Moore: ciclo di ricerca}
\begin{algorithmic}[1]
\Require testo $T$ di lunghezza $m$, pattern $P$ di lunghezza $n$
\Ensure le posizioni iniziali di tutte le occorrenze di $P$ in $T$
\State calcola $R_j(x)$, $L'$ e $\ell'$ con il preprocessing di $P$
\State $k \gets n$ \Comment{$P[n]$ è allineato con $T[k]$}
\While{$k \le m$}
  \State $i \gets n+1$; \quad $h \gets k+1$
  \While{$i>1$ \textbf{ e } $P[i-1]=T[h-1]$}
    \State $i \gets i-1$; \quad $h \gets h-1$
  \EndWhile
  \If{$i=1$}
    \State riporta l'occorrenza che inizia in $k-n+1$
    \State $k \gets k+\sigma(1)$
  \Else
    \State $x \gets T[h-1]$ \Comment{il carattere «cattivo»}
    \State $k \gets k+\max\set{\,(i-1)-R_{i-1}(x),\ \sigma(i)\,}$
  \EndIf
\EndWhile
\end{algorithmic}
\end{algorithm}

L'invariante del ciclo interno è $\sub{P}{i}{n}=\sub{T}{h}{k}$; all'uscita, o $i=1$
e si è trovata un'occorrenza, oppure $P[i-1]\neq T[h-1]$ e si è nella situazione
descritta nella \autoref{sec:bm-regole}.

\begin{table}[H]
\centering
\begin{tabular}{ll}
\toprule
\textbf{Fase} & \textbf{Costo} \\
\midrule
Preprocessing BCR, matrice & $\Th(n\abs{\alfa})$ tempo e spazio, interrogazioni $\Oh(1)$ \\
Preprocessing BCR, liste   & $\Th(n+\abs{\alfa})$ spazio, costo ammortizzato equivalente \\
Preprocessing GSR ($Z$ su $P^r$) & $\Th(n)$ tempo e spazio \\
Ricerca, caso migliore     & $\Th(m/n)$ caratteri ispezionati (sublineare) \\
Ricerca, caso peggiore     & $\Th(mn)$ (quadratico) \\
\bottomrule
\end{tabular}
\caption{Riepilogo dei costi di Boyer--Moore.}
\end{table}

Il caso peggiore quadratico è quello già visto di $T=a^{m}$, $P=a^{n}$: è proprio
per eliminarlo, senza rinunciare alle due euristiche, che serve la variante di
Apostolico--Giancarlo.
